\section{Kết luận}
Nghiên cứu đề xuất GeoFormerDock, một kiến trúc học sâu đa nhiệm gọn nhẹ (1.59 triệu tham số) cho bài toán phân loại tư thế gắn kết và dự đoán ái lực protein--phối tử. Trên khung đánh giá chung xây dựng cho bốn kiến trúc, GeoFormerDock đạt vị trí không bị trội trên đường biên Pareto giữa hai nhiệm vụ, dẫn đầu về Balanced Accuracy và PR-AUC ở nhánh phân loại tư thế với quy mô tham số nhỏ hơn đáng kể so với các kiến trúc CNN chuyên biệt, trong khi vẫn còn khoảng cách có ý nghĩa thống kê với GNINA Dense ở nhánh ái lực (riêng với GNINA Default2018 và Pafnucy, khác biệt chỉ rõ ràng ở MAE/RMSE).

Một số điểm cần lưu ý khi diễn giải các kết quả trên: (i) checkpoint tốt nhất và điểm dừng sớm hiện được chọn trên chính tập kiểm tra do chưa có tập validation độc lập; giữa hai epoch liên tiếp (98 và 100), Balanced Accuracy đo được dao động tới $0.163$, một mức chênh lệch đáng cân nhắc khi diễn giải kết quả; (ii) các số liệu hiện dựa trên một lần chia train/test, chưa có điều kiện chạy lặp lại hay chia theo cluster protein/pocket để đánh giá đầy đủ hơn khả năng tổng quát hóa; (iii) nhóm mô hình đồ thị (PotentialNet, EquiBind, TankBind) không được đưa vào so sánh vì hoạt động trên biểu diễn đồ thị nguyên tử, không tương thích với khung voxel thống nhất của nghiên cứu này; so sánh chéo giữa hai họ biểu diễn đòi hỏi thiết kế thực nghiệm riêng; (iv) các chỉ số ái lực được tính trên 623/4.618 mẫu có nhãn ái lực dương; (v) việc tách bạch đóng góp của kiến trúc so với các yếu tố huấn luyện (hàm mất mát, cân bằng lớp ở mức minibatch) vẫn là hướng cần làm rõ thêm.
